\section{Landmark physics rewrite plan}
\label{sec:landmark-physics-rewrite-plan}

\subsection{Claim attacked}

\textbf{Claim attacked.}  The integrated manuscript now satisfies the
physical-first principle: charge/BPS/duality data precede the
lattice/modular formalism, every physical assertion has the right status,
and the construction is derived from the physical operation rather than
from a later automorphic convention.

\textbf{Attack result.}  False as written.  The integration has healed the
largest overclaim: the Borcherds algebra is now a denominator algebra, not
the microscopic BPS-state algebra.  The remaining violations are more
structural.  The paper still sometimes names the universal
Gram-invariant index lattice as the physical charge lattice; it still
lets automorphic covariance appear before the protected-index source; and
it still says in a few places that the physical operation fixes an object
which is fixed only after adding an explicit physical dictionary.

\subsection{Manuscript anchors}

\begin{center}
\small
\begin{tabular}{p{0.18\textwidth}|p{0.37\textwidth}|p{0.35\textwidth}}
\textbf{Anchor in current \texttt{proj.tex}}&
\textbf{Violation}&
\textbf{Required upgrade}\\ \hline
55--115, abstract&
The abstract says the square root is ``singled out by the physics'' and
then gives the automorphic denominator algebra.  This compresses three
different levels: OP scalar square, virtual K3 half-index determinant,
and denominator algebra.&
Replace by a status-separated abstract: OP scalar square is proved;
the Borcherds determinant is the normalized automorphic square root;
the microscopic BPS factorization realization is open.\\ \hline
82--88, abstract&
``The BPS charge pairing'' and ``same charge lattice'' are too strong
before the physical charge lattice has been constructed.  What is proved
in the paper is the Gram-invariant index lattice and its anti-isometry to
\(\Lambda^{2,1}_{II}\).&
Use ``Gram-invariant index lattice'' until the physical realization
assumption is stated.  Reserve ``physical charge lattice'' for the
microscopic compactification lattice and its projection to
\((Q^2/2,Q\cdot P,P^2/2)\).\\ \hline
121--178, duality principle&
The section begins with an automorphic section and only then introduces
the charge matrix.  This is backwards for a physical-first paper.&
Begin with the physical two-charge input as an assumption, derive
\(T(n,l,m)\), \(4nm-l^2\), and the trace character, then state the
automorphic section as the duality-covariant completion.\\ \hline
167--178, charge matrix proof&
``Thus the variables in the Borcherds product are the protected BPS
charge invariants'' is stronger than the proof.  The proof shows this
only after assuming the protected index factors through the two-charge
Gram matrix.&
Downgrade to: ``Under the Gram-invariant protected-index hypothesis,
these are the variables of the BPS index.''\\ \hline
195--204, cusp polarization&
``It gives exactly the effective semigroup'' can read as a physical
wall statement.  It is a Fourier/Fock polarization at the standard cusp.&
State it as a convention fixed by the ordered cusp limit.  Its physical
wall-crossing meaning is not used unless a microscopic stability
condition is supplied.\\ \hline
217--223, one-particle index&
``The charge lattice is \(\Gamma_{\mathrm{BPS}}=\mathbb Z^3\)''
overidentifies the invariant triple lattice with the compactification
charge lattice.&
Write \(\Gamma_{\mathrm{gram}}=\mathbb Z^3\) or
\(\Gamma_{\mathrm{idx}}=\mathbb Z^3\).  Then state the map from physical
charges to it as an assumption or dictionary.\\ \hline
551--668, protected D-brane index and scalar square&
This is the real physical scalar-square source, but it appears after the
K3 half-index determinant and after the initial duality section.&
Move this material before the determinant construction.  The logical
entrance should be: protected D-brane index, OP theorem, DT/protected
dictionary status, then square-root determinant.\\ \hline
839--861, automorphic duality datum&
``The geometric duality group is \(\PSp_4(\mathbb Z)\)'' is stated as
if it has already been derived from the compactification.&
Rename to ``automorphic duality group used here''.  Add: its
identification with the physical duality group is physical input, not a
consequence of the exterior-square lattice computation.\\ \hline
865--869 and 1030--1074, Weyl chamber&
``same BPS charge lattice'' still appears in the denominator
polarization.  The healed statement is that the Weyl chamber is the
denominator polarization of the same Gram-invariant index lattice,
restricted to active Jacobi support.&
Replace ``BPS charge lattice'' by ``Gram-invariant index lattice'' here.
Keep the active-support lemma; it is the correct repair.\\ \hline
1253--1277, denominator algebra square root&
This is mostly healed.  The title still tempts the reader to hear a
physical bracket-level square root.&
Keep the theorem but rename internally to ``Denominator-algebra square
root''.  In prose say ``algebraic denominator square root'', never
``physical BPS square root''.\\ \hline
1372--1455, Hall--Drinfeld criterion&
The criterion correctly lists missing structures, but condition (1)
again says the protected charge lattice is \(\mathbb Z^3\), and line
1448 says OP fixes the protected character.&
Change condition (1) to a projection to the Gram-invariant index lattice.
Change ``OP fixes the protected character'' to ``OP fixes the primitive
reduced curve-counting character; it becomes protected BPS character only
after the protected-index dictionary.''\\ \hline
1679--1691, index equality&
``one-particle BPS index coefficient'' can still be heard as an honest
state-space coefficient.&
Use ``virtual K3 half-index coefficient'' or ``protected Euler
coefficient''.  The following sentence may then say that a physical
one-particle interpretation is additional.\\ \hline
3231--3358, Vol.~III row material&
This material is well quarantined but belongs after all \(N=1\) physics
and denominator results.  Its physical row language should not leak into
the main theorem.&
Keep as appendix-only compatibility and obstruction material.  Do not
use it to motivate the \(N=1\) construction before the actual
charge/BPS derivation.
\end{tabular}
\end{center}

\subsection{Replacement theorem order}

The global rewrite should use this theorem order.

\begin{enumerate}
\item \textbf{Physical scalar problem.}
State \(X=S\times E\), the D6--D2--D0 charge notation, the protected
index \(\Omega_X\), and the status of the DT/protected dictionary.  Then
state Oberdieck--Pixton as the unconditional primitive reduced
curve-counting theorem:
$$
Z^X_{\mathrm{OP}}=-4096\,\Delta_5^{-2}.
$$
Status: proved for OP; physical DT/BPS interpretation is conjectural or
assumptive as already recorded.

\item \textbf{K3 protected one-particle input.}
State the Ramond-sector index
$$
Z_{\mathrm{K3}}=2\phi_{0,1}
$$
and the right-moving localization lemma.  Then define the virtual
half-index class \(\mathbb U_{n,l}\).  Status: protected index is
physical; the half-index object is virtual and noncanonical.

\item \textbf{Gram-invariant index lattice.}
Assume the protected sector is read through two-charge Gram invariants:
$$
n=Q^2/2,\qquad l=Q\cdot P,\qquad m=P^2/2.
$$
Define \(\Gamma_{\mathrm{gram}}=\mathbb Z^3\) and
$$
\langle\gamma,\gamma'\rangle_{\mathrm{gram}}
=2(nm'+n'm)-ll'.
$$
Status: mathematical consequence after the physical Gram-invariance
input.  Do not call this the full physical charge lattice.

\item \textbf{Cusp polarization and active support.}
Define \(\Gamma_{\mathrm{eff}}\) from the ordered cusp limit, then define
\(\Gamma_{\mathrm{act}}\) by \(f(nm,l)\ne0\).  Prove the active-support
cone before introducing the Weyl chamber as a chamber:
$$
\alpha(\Gamma_{\mathrm{act}})
\subset
\mathbb Z_{\ge0}\delta_1+\mathbb Z_{\ge0}\delta_2+\mathbb Z_{\ge0}\delta_3,
$$
and the three simple rays are active.

\item \textbf{Virtual Fock determinant.}
Build \(\mathcal V\) from \(\mathbb U_{nm,l}\), state the Fock
superdeterminant lemma, and prove
$$
\mathcal D_X
=64q^{1/2}r^{1/2}s^{1/2}
\prod_{\Gamma_{\mathrm{eff}}}(1-q^nr^ls^m)^{f(nm,l)}.
$$
Then invoke Borcherds--Gritsenko--Nikulin to identify
\(\mathcal D_X=\Delta_5\).  Status: determinant is virtual; equality to
\(\Delta_5\) is the Borcherds product theorem.

\item \textbf{Automorphic completion.}
Only after the determinant is constructed, introduce the exterior-square
model:
$$
\PSp_4(\mathbb Z)\simeq \SO_+(\Lambda^{3,2}).
$$
Status: theorem about lattices and period domains.  Its use as physical
duality is an explicit physical assumption.

\item \textbf{Weyl chamber and denominator cone.}
Use the anti-isometry
$$
\alpha(n,l,m)=2nf_2-lf_3+2mf_{-2}
$$
to identify the active-support cone with the type-II simple-root cone.
Then compute the Weyl vector and multiplier restriction.  Status:
mathematical denominator polarization of the Gram-invariant index
lattice.

\item \textbf{Denominator algebra.}
Construct \(\mathcal A_{\mathrm{den}}=\mathfrak g_{\Delta_5}\).  Prove
the denominator identity and the signed superdimension equality.  Status:
proved denominator algebra, not microscopic BPS algebra.

\item \textbf{Formal current envelope.}
Construct
$$
\mathsf{Den}_{\Delta_5,E}
=U_E^{\mathrm{ch}}(\mathcal A_{\mathrm{den}}\otimes\mathcal O_E).
$$
Status: formal BD current envelope.  Not a compact \(K3\times E\)
factorization algebra.

\item \textbf{Microscopic realization criterion.}
Move the Hall--Drinfeld criterion after the formal envelope.  Make it the
only place where a physical bracket-level realization is discussed.
Status: open until the compact Hall cosheaf, protected integration,
orientation character, and Borcherds relations inside the Hall bracket
are constructed.

\item \textbf{Catalogue and Vol.~III row map.}
Keep the Clery--Gritsenko and Vol.~III row material in an appendix or
late section.  It must not supply an early physical motivation for the
\(N=1\) theorem.
\end{enumerate}

\subsection{Sample replacements}

\paragraph{Opening paragraph.}
The scalar input is a protected enumerative partition function.  For
\(X=S\times E\), Oberdieck--Pixton prove the primitive reduced
curve-counting identity
$$
Z^X_{\mathrm{OP}}=-4096\,\Delta_5^{-2}.
$$
The Hilbert-scheme DT and microscopic BPS readings require additional
dictionary data.  The paper constructs the canonical automorphic square
root on the unconditional OP branch and records exactly which further
physical data would identify it with a compact BPS factorization
algebra.

\paragraph{Square-root paragraph.}
The square root used below is not an arbitrary analytic branch of
\(\Delta_5^2\).  It is the normalized Borcherds determinant
$$
\mathcal D_X=\Delta_5,
$$
whose product expansion is built from the virtual half of the K3
Ramond-sector index.  This is a protected-index construction.  It does
not construct a canonical half of the K3 Hilbert space and does not
construct a microscopic \(K3\times E\) BPS algebra.

\paragraph{Charge-lattice paragraph.}
Let \(\Gamma_{\mathrm{gram}}=\mathbb Z^3\) be the lattice of
rank-two Gram invariants
$$
\gamma=(n,l,m)=\left(Q^2/2,\ Q\cdot P,\ P^2/2\right).
$$
The bilinear form
$$
\langle\gamma,\gamma'\rangle_{\mathrm{gram}}
=2(nm'+n'm)-ll'
$$
is the polarization of the discriminant \(4nm-l^2\).  The map
$$
\alpha(n,l,m)=2nf_2-lf_3+2mf_{-2}
$$
is an anti-isometry from this Gram-invariant index lattice to
\(\Lambda^{2,1}_{II}\).  The assertion that every relevant protected
BPS index of the compactification factors through this lattice is a
physical input, not a conclusion of the lattice computation.

\paragraph{Automorphic duality paragraph.}
The exterior-square construction proves the lattice isomorphism
$$
\PSp_4(\mathbb Z)\simeq\SO_+(\Lambda^{3,2}).
$$
This is the automorphic duality group of the Siegel period domain used by
the Igusa form.  Identifying it with the physical duality group acting on
the protected compactification sector is a separate physical
assumption.  The subsequent modular covariance statements use only this
automorphic action and the stated assumption.

\paragraph{Hall realization paragraph.}
The Oberdieck--Pixton scalar square fixes a primitive reduced
curve-counting character.  Under the protected-index dictionary it is
expected to be the protected BPS character.  It does not construct the
compact oriented critical Hall cosheaf, the protected integration map,
the Hall structure constants, the orientation character
\(\nu_{\Delta_5}\), or the comparison with the Borcherds relations.
Those data are precisely the microscopic realization problem.

\subsection{Broken statements to upgrade}

\begin{enumerate}
\item Replace every main-text occurrence of
``the charge lattice is \(\Gamma_{\mathrm{BPS}}=\mathbb Z^3\)''
by ``the Gram-invariant index lattice is
\(\Gamma_{\mathrm{gram}}=\mathbb Z^3\)'', unless the sentence is inside a
conditional physical-realization hypothesis.

\item Replace ``the geometric duality group is
\(\PSp_4(\mathbb Z)\)'' by ``the automorphic duality group used here is
\(\PSp_4(\mathbb Z)\); its physical identification is input''.

\item Replace ``the square root singled out by the physics'' by
``the normalized Borcherds determinant selected by the protected-index
normalization''.  Add the physical square-root interpretation only under
the protected-index dictionary.

\item Replace ``OP fixes the protected character'' by ``OP fixes the
primitive reduced curve-counting character; it becomes protected BPS
character under the protected-index dictionary''.

\item Replace ``same BPS charge lattice'' in the Weyl section by
``same Gram-invariant index lattice''.  Then cite the active-support
lemma as the reason the denominator cone is the simple-root cone.

\item Replace ``one-particle BPS index coefficient'' in the determinant
forgetfulness theorem by ``virtual K3 half-index coefficient'' unless an
honest one-particle Hilbert lift has been assumed.

\item Move all Vol.~III row-map and boundary-row material after the
\(N=1\) theorem.  It is compatibility and obstruction material, not a
premise of the Igusa square-root construction.
\end{enumerate}

\subsection{Files changed}

Only this owned file was changed:
$$
\texttt{agent\_material/08\_landmark\_physics\_rewrite\_plan.tex}.
$$

\subsection{Commands run}

\begin{verbatim}
sed -n '1,240p' AGENTS.md
sed -n '1,220p' CLAUDE.md
sed -n '128,148p' /Users/raeez/ecosystem/INVARIANTS.md
sed -n '133,148p' /Users/raeez/ecosystem/AGENTS-HARNESS.md
git status --short
rg -n "^\\section|^\\subsection|physical|BPS|charge|duality|..." proj.tex
nl -ba proj.tex | sed -n '50,122p'
nl -ba proj.tex | sed -n '121,214p'
nl -ba proj.tex | sed -n '215,405p'
nl -ba proj.tex | sed -n '542,736p'
nl -ba proj.tex | sed -n '739,862p'
nl -ba proj.tex | sed -n '863,1125p'
nl -ba proj.tex | sed -n '1210,1470p'
nl -ba proj.tex | sed -n '1658,1732p'
nl -ba proj.tex | sed -n '1900,2018p'
nl -ba proj.tex | sed -n '2570,2768p'
nl -ba proj.tex | sed -n '3230,3368p'
rg -n "bracket|BKM|Igusa|Delta_5|physical|open|Hall|Drinfeld|..." \
  /Users/raeez/calabi-yau-quantum-groups/FRONTIER.md
sed -n '1,320p' agent_material/07_chiral_koszul_factorization_boundary.tex
sed -n '1,260p' agent_material/07_nonabelian_promotion_f5_f8.tex
rg -n "charge lattice is|same BPS charge lattice|physical duality|..." \
  proj.tex agent_material -S
\end{verbatim}

\subsection{Remaining open obligations}

\begin{enumerate}
\item Construct the compactification charge lattice and prove the
projection to \((Q^2/2,Q\cdot P,P^2/2)\) is the complete invariant of
the protected sector used by the product.

\item Prove that the physical duality group acting on the protected
\(K3\times E\) sector is the automorphic group
\(\PSp_4(\mathbb Z)\) used in the exterior-square model, or state the
precise subgroup/coset actually supplied by the compactification.

\item Prove the protected-index dictionary
$$
\Omega_X(h,d,n)=\mathbf{DT}^X_{n,(\beta_h,d)}
$$
in the normalization and chamber of the paper, or keep it as a
conjecture.

\item Construct an honest chiral one-particle Hilbert lift of the
virtual half-index, or keep the determinant entirely in
\(K_0\)-language.

\item Construct the compact oriented critical Hall cosheaf for
\(K3\times E\), its protected integration map, its Hall
super-commutator, and the comparison of its primitive bracket with
\(\mathfrak g_{\Delta_5}\).

\item Prove that the automorphic real-root hyperplanes are microscopic
walls of marginal stability for the compactification.  Until then, they
are denominator walls.

\item Build the Vol.~III row map
\(\Phi^{\mathrm{V3}}_j\leftrightarrow F_i\) before using any boundary
row as a physical partition-function theorem.
\end{enumerate}
